\section{Organización del código fuente}
\label{sec:codebase}

El repositorio refleja directamente la separación entre los dos modos
de despliegue descritos en las secciones anteriores:

\begin{figure*}[!t]
\centering
\begin{verbatim}
PARTS_SORTING_PLATFORM/
|-- shared/            clasificador por espectro (FFT), compartido entre modos
|-- mobile-app/        Expo + React Native + TypeScript
|   `-- src/
|       |-- api/
|       |-- components/
|       |-- context/
|       `-- screens/
|-- CLOUD_MODE/
|   |-- auth-service/
|   |-- api-gateway/
|   |-- classification-service/
|   |-- mqtt-broker/
|   |-- firmware/esp32s3_result_display/
|   `-- scripts/
|-- LOCAL_MODE/
|   |-- local-server/
|   |   |-- app/         FastAPI: clasificador YOLO, medicion de diametro
|   |   |-- training/    dataset + entrenamiento YOLOv8
|   |   |-- static/      paginas servidas al telefono
|   |   `-- tests/
|   |-- firmware/esp32_local_http/
|   |   `-- sprits/      sketches de prueba por modulo
|   `-- scripts/
|-- docs/
`-- REPORT/
\end{verbatim}
\caption{Estructura de directorios de \texttt{PARTS\_SORTING\_PLATFORM}.}
\label{fig:codebase_tree}
\end{figure*}

El único módulo verdaderamente compartido entre ambos modos de
despliegue es \texttt{shared/spectrum\_classifier.py}. En
\texttt{CLOUD\_MODE} este módulo sigue siendo el clasificador
efectivamente usado por \textit{classification-service}; en
\texttt{LOCAL\_MODE}, en cambio, la ruta de clasificación en vivo fue
migrada a un modelo YOLOv8 entrenado específicamente sobre el banco de
pruebas (Sección~\ref{sec:local_mode}), de modo que el clasificador por
espectro quedó presente en el código pero sin uso activo en ese modo.

\subsection{Pertinencia del clasificador por espectro frente a las limitaciones observadas de YOLO}

Durante las pruebas del modo local se observó que el modelo YOLO,
entrenado sobre un conjunto de datos todavía pequeño y desbalanceado
por clase, puede ser inconsistente entre fotografías consecutivas de
la misma pieza física: la clasificación alternó entre dos clases
distintas y la estimación de diámetro varió apreciablemente sin que la
pieza se moviera, pese a que el modelo reportaba una confianza muy alta
(superior al 95\,\%) en ambos casos. Esto es un síntoma típico de
sobreajuste con pocos datos de entrenamiento, no un error de
implementación.

En ese contexto, el clasificador por espectro conserva un valor
práctico concreto: al ser un método completamente distinto (basado en
bandas de energía radiales de la transformada de Fourier, sin red
neuronal), no comparte los sesgos de un modelo entrenado con pocos
ejemplos, y podría usarse como una segunda opinión independiente o como
mecanismo de verificación cruzada frente a una predicción de YOLO de
alta confianza pero potencialmente errónea. Esta combinación
(\textit{ensemble}) no está implementada en el código actual; se deja
planteada como mejora futura concreta, no como una capacidad ya
existente.

\subsection{Una condición de carrera real, y cómo se resolvió}

Es importante distinguir lo anterior de una condición de carrera real
que sí se encontró y corrigió en el código del modo local: al pasar de
una señal de pausa de disparo único a una señal de pausa por nivel
(que permanece activa hasta que se indica explícitamente reanudar), la
página que captura la fotografía consultaba ese estado repetidamente y
disparaba una captura nueva en cada consulta mientras la pausa seguía
activa, generando decenas de fotografías para una sola detención de la
banda. La corrección fue estrictamente de control de estado —una
bandera que recuerda si ya se capturó una fotografía para la pausa
actual— y no tiene relación alguna con el procesamiento espectral de
la señal.

\subsection{Estado de la aplicación móvil}

El código de la aplicación móvil está completo: pantallas de
autenticación, captura de fotografía y consulta de historial, cliente
HTTP con renovación automática de token, y almacenamiento seguro de
credenciales. Se verificó de forma estática (compilación de tipos sin
errores y verificación de configuración del proyecto), pero, dada la
complejidad de disponer de un emulador o dispositivo físico configurado
en este entorno de desarrollo, no se llegó a ejecutar en vivo contra un
backend real. El código existe y está listo para esa prueba; lo que
falta es la verificación funcional end-to-end, no su implementación.
